\subsection{Data sample and event selection\label{sec:data}}

\subsubsection{\PbPb\ Data\label{sec:data_PbPb}}

The \PbPb\ data used in this analysis were recorded during the 2015 and 2018 \PbPb\ run 
  at $\sqrtsnn=5.02$~\TeV\ in the ATLAS experiment. 
The ATLAS experiment and the trigger system are decribed 
  in Refs.~\cite{PERF-2007-01}--\cite{TRIG-2016-01}.
The 2015 data were recorded over 33 separate runs and the 2018 data were 
  recorded over 39 separate runs. 
After applying a selection of runs where the detector was fully operational, 
  a total integrated luminosity of $\sim$1.94~\inb was used for this analysis, 
  of which 0.50~\inb~\cite{Arratia:2152730} was taken in 2015
  and 1.44~\inb\ in 2018.


Events were selected based using the following di-muon trigger 
\begin{enumerate}
\item \verb=HLT_2mu4=. A level 1 (L1) - this trigger requires the \verb=L1_2MU4= 
     which selects events with two muons above the \verb=L1_MU4= threshold. 
   A High Level Trigger (HLT) muon reconstruction algorithm is then applied to the regions of 
     interest identified at L1 which used information from the inner detector.
   Finally, both HLT-reconstructed muons are required to have a transverse momentum above 4~\GeV.\
   This trigger was unprescaled during the 2015 and 2018 run periods.

\item \verb=HLT_mu4_mu4noL1=. At L1,  this trigger requires the \verb=L1_MU4=
  trigger which selects events with a single muon above the \verb=L1_MU4= 
    threshold. 
  A HLT muon reconstruction algorithm is then applied to the 
    region of interest identified at L1 for the first muon and a full scan
    is performed for the second muon.
  Both HLT-reconstructed muons are required to have a transverse momentum above 4~\GeV.\
  This trigger was prescaled for some small parts of the 2015 and 2018 \PbPb\ runs.
  However, the HLT hypos for the \verb=HLT_mu4= and \verb=HLT_mu4noL1= 
    chains are different, resulting in \verb=HLT_mu4_mu4noL1= not being 
    a superset of the \verb=HLT_2mu4= trigger.
  Additionally, because the \verb=HLT_mu4_mu4noL1= was prescaled in a few runs,
    while the \verb=HLT_2mu4= trigger was not, including the \verb=HLT_2mu4=
    trigger gains some additional statistics.
\end{enumerate}


Additional minimum-bias and single-muon triggered samples were used to study the trigger efficiency. 
The minimum-bias data set was taken using a combination of the following
  mutually exclusive total transverse energy and ZDC triggers. 
\begin{enumerate}
  \item \verb=HLT_noalg_cc_L1TE600_0ETA49= : This required at least 600~\GeV\ 
    of total transverse energy (\ET) in the L1 calorimeter system. 
  \item \verb=HLT_noalg_pc_L1TE50_VTE600_0ETA49= : This trigger required that the 
    total \ET\ in the L1 calorimeter system be between 50 and 600~\GeV.
  \item \verb=HLT_mb_sptrk_ion_L1ZDC_A_C_VTE50= : This trigger selected events that 
    had an L1 \ET\ less than 50~\GeV\ but satisfied a ZDC coincidence trigger. 
    At the HLT, this trigger required at least one track reconstructed in the inner detector 
      in order to reject empty events.
\end{enumerate}
The single-muon triggered sample was taken using a combination of the following triggers:
\begin{itemize}
\item \verb=HLT_mu4=
\item \verb=HLT_mu6=
\item \verb=HLT_mu8=.
\end{itemize}

Examples of the trigger menus for the 2015 and 2018 runs can be found
  from the run query page which are linked here for the 
  \href{https://atlas-runquery.cern.ch/query.py?q=find+run+286711++%2F+show+all}{2015 Run}
  and
  \href{https://atlas-runquery.cern.ch/query.py?q=find+run+365512+%2F+show+all}{2018 Run}.
The following GRLs were used for the 2015 and 2018 data:
\begin{itemize}
\item \verb=data18_hi.periodAllYear_DetStatus-v106-pro22-14_Unknown_PHYS= \newline 
      \verb=_HeavyIonP_All_Good.xml=
\item \verb=data15_hi.periodAllYear_DetStatus-v75-repro20-01_DQDefects= \newline 
      \verb=-00-02-02_PHYS_HeavyIonP_All_Good.xml=
\end{itemize}
The 2018 GRL excludes runs 367318-367384 which had issues with the toroid magnet. 
In addition to the GRL requirements, in the data analysis, events were required 
  to have a primary reconstructed vertex and contain no detector or data 
  acquisition errors (e.g. events corrupted from TTC restarts).
The following datasets (or dataset containers) were used in the Analysis:
\begin{itemize}
\item 2018 Data: {\small \verb=data18_hi:data18_hi.<RunNo>.physics_HardProbes.merge.AOD.f1021_m2037=}
\item 2018 Data: {\small \verb=data18_hi:data18_hi.<RunNo>.physics_HardProbes.merge.AOD.f1021_m2037=}
\item 2018 Data: {\small \verb=data18_hi:data18_hi.<RunNo>.physics_HardProbes.merge.AOD.f1030_m2048=}
\item 2015 Data: {\small \verb=data15_hi.periodL.physics_HardProbes.PhysCont.AOD.pro22_v01=}
\end{itemize}


The centrality of the \PbPb\ collisions is characterised by the total
  transverse energy measured in the FCal modules, \ETfcal. 
The \ETfcal\ distribution obtained in minimum-bias collisions is partitioned 
  into separate ranges of \ETfcal\ referred to as centrality
  classes~\cite{Miller:2007ri,Alver:2008aq,HION-2011-05}. 
Each class is defined by the fraction of the distribution contained by 
  the interval, e.g. the 0--10\% centrality class, which corresponds 
  to the most central collisions, contains the 10\% of minimum-bias (i.e. inclusive)
  events with the largest \ETfcal.
An average number of participants, $\langle N_{\mathrm{part}} \rangle$, 
  is defined for each centrality class using a standard Glauber MC 
  procedure~\cite{Perepelitsa:2212936,CentralityRec}. 
The Heavy-ion group provides standard centrality bins 
  for Pb+Pb events based on the $\ETfcal$, that covers the 0--80\%  
  centrality-percentile range, with the 80\% centrality value corresponding to 
  an \ETfcal\ value of 63.7~\GeV.
Beyond this range it is difficult to assign a centrality-percentile classification 
  to an event.
Figure~\ref{fig:events_in_triggers} shows the number of events
  and the centrality distribution for events collected by the 
  logical OR of the two di-muon triggers.

\begin{figure}
\centering
\includegraphics[width=0.5\textwidth]%
{figures/AllFigs/can_cent_pTBar_GT5_EffCor}
\caption{
The number of events selected by the di-muon triggers 
  as a function of centrality.
%The number of events selected by the logical OR of 
%  the HLT\_2mu4 and HLT\_mu4\_mu4noL1 triggers, 
%  as a function of centrality.
%The distribution is shown only over the 0--80\% centrality interval (see text).
}
\label{fig:events_in_triggers}
\end{figure}



\subsubsection{\pp\ Data\label{sec:data_pp}}
The \pp\ data at \sqs=5.02~\TeV\ were recorded during a dedicated 
  reference run in 2017.
The average interaction per bunch-crossing ($\langle\mu\rangle$) 
  for this period varied between 0.4 and 4.
The data used here were recorded by the \verb=HLT_2mu4= trigger.
This trigger was active during the whole data taking period corresponding
  to a total integrated luminosity of 0.26~\ifb.
The following GRL and dataset containers were used for the 2017 \pp\ data:
\begin{itemize}
\item GRL: \verb=data17_5TeV.periodAllYear_DetStatus-v98-pro21-16_Unknown_PHYS= \newline 
           \verb=_StandardGRL_All_Good_25ns_ignore_GLOBAL_LOWMU.xml=
\item Dataset: \verb=data17_5TeV.periodM.physics_Main.PhysCont.AOD.pro23_v02=
\end{itemize}



\subsection{Monte Carlo samples\label{sec:mc_samples}}
The MC samples are used for two purposes: 
  1) to obtain the muon reconstruction efficiencies, and 
  2) to obtain the momentum-imbalance dsitributions for
     signal and for background muons.
The \textit{signal} muons are muons that arise
  from the decay of HF particles.
The background muons are muons that arise from secondary 
  interactions in the calorimeter, or from decays from pions and Kaons
  and from combinatoric background that arise from
  mis-associations of MS tracks with ID tracks.
The momentum-imbalance distributions obtained from the MC are used to estimate
  and correct for the contribution of background muons in data analysis.
The following MC samples are used in this analysis: 

\subsubsection{\PbPb\ MC Samples~\label{sec:PbPbMc}}
The \PbPb\ muon reconstruction efficiencies and signal-templates  
  for momentum-imbalance distributions (for true muons) are obtained 
  from HIJING-overlay Monte Carlo (MC) samples. 
These samples were generated for the analysis ANA-HION-2020-10, 
  which studies the $\gamma\gamma\rightarrow \mumu$ in \PbPb\ events,
  but can also be used in the present analysis.
The details for these MC samples can be found at the following JIRAs
\begin{itemize}
\item \textbf{HIJING Overlay 2015 conditions: } ATLHI-304~\cite{ATLHI-304}
\item \textbf{HIJING Overlay 2018 conditions: } ATLHI-313~\cite{ATLHI-313}
\end{itemize}
The \starlight\ event generator~\cite{Klein:2016yzr} was used to generate
  $\gamma\gamma\rightarrow \mumu$ events. 
To account for the effects of the \PbPb\ underlying event, an \emph{overlay}
  procedure was utilized; simulated events $\gamma\gamma\rightarrow \mumu$ 
  were overlaid directly onto MC events generated by the HIJING event generator.
Through this procedure, the effects of the underlying event (UE) 
  on the muon measurement in the overlay sample closely match 
  those present in the data.

The following HIJING overlay Pythia samples were used to generate 
  momentum-imbalance templates for the background muons (i.e. misidentified muons):
\begin{itemize}
  \item {\small \verb=mc16_valid.420011.Pythia8EvtGen_A14NNPDF23LO_jetjet_JZ1R04.recon.AOD.e4108_s2873_r10756=}
  \item {\small \verb=mc16_valid.420012.Pythia8EvtGen_A14NNPDF23LO_jetjet_JZ2R04.recon.AOD.e4108_s2873_r10756=}
  \item {\small \verb=mc16_valid.420013.Pythia8EvtGen_A14NNPDF23LO_jetjet_JZ3R04.recon.AOD.e4108_s2873_r10756=}
  \item {\small \verb=mc16_valid.420014.Pythia8EvtGen_A14NNPDF23LO_jetjet_JZ4R04.recon.AOD.e4108_s2873_r10756=}
\end {itemize}


\subsubsection{\pp\ MC Samples}
For the \pp\ analysis, the reconstruction efficiency is directly 
  obtained from the MCP recommendations in Ref.~\cite{Rel21Uncertainties}.
The MC samples are used to only obtain the momentum-imbalance 
  distributions for signal and background muons.
The following Pythia8 dijets samples (2017 Conditions) were used 
  (\href{https://twiki.cern.ch/twiki/bin/viewauth/AtlasProtected/HIJetMCSamples#Pythia8_dijets_8M_per_sample_in}{link to twiki}): 
\begin{itemize}
  \item {\small \verb=mc16_5TeV.420010.Pythia8EvtGen_A14NNPDF23LO_jetjet_JZ0R04.recon.AOD.e4108_s3238_r11199=}
  \item {\small \verb=mc16_5TeV.420011.Pythia8EvtGen_A14NNPDF23LO_jetjet_JZ1R04.recon.AOD.e4108_s3238_r11199=}
  \item {\small \verb=mc16_5TeV.420011.Pythia8EvtGen_A14NNPDF23LO_jetjet_JZ1R04.recon.AOD.e6608_s3238_r11199=}
  \item {\small \verb=mc16_5TeV.420012.Pythia8EvtGen_A14NNPDF23LO_jetjet_JZ2R04.recon.AOD.e4108_s3238_r11199=}
  \item {\small \verb=mc16_5TeV.420012.Pythia8EvtGen_A14NNPDF23LO_jetjet_JZ2R04.recon.AOD.e6608_s3238_r11199=}
  \item {\small \verb=mc16_5TeV.420013.Pythia8EvtGen_A14NNPDF23LO_jetjet_JZ3R04.recon.AOD.e4108_s3238_r11199=}
  \item {\small \verb=mc16_5TeV.420013.Pythia8EvtGen_A14NNPDF23LO_jetjet_JZ3R04.recon.AOD.e6608_s3238_r11199=}
  \item {\small \verb=mc16_5TeV.420014.Pythia8EvtGen_A14NNPDF23LO_jetjet_JZ4R04.recon.AOD.e4108_s3238_r11199=}
  \item {\small \verb=mc16_5TeV.420014.Pythia8EvtGen_A14NNPDF23LO_jetjet_JZ4R04.recon.AOD.e6608_s3238_r11199=}
  \item {\small \verb=mc16_5TeV.420015.Pythia8EvtGen_A14NNPDF23LO_jetjet_JZ5R04.recon.AOD.e4108_s3238_r11199=}
  \item {\small \verb=mc16_5TeV.420015.Pythia8EvtGen_A14NNPDF23LO_jetjet_JZ5R04.recon.AOD.e6608_s3238_r11199=}
\end{itemize}
